\documentclass[conference]{IEEEtran}
\usepackage{amsmath}
\usepackage{booktabs}
\usepackage{graphicx}

\title{PATV-X: Training-Free Explainable Deepfake Detection via Multi-Scale Frequency Forensics}

\author{\IEEEauthorblockN{George David Tsitlauri}
\IEEEauthorblockA{
gdtsitlauri@gmail.com\\
gdtsitlauri.dev\\
github.com/gdtsitlauri
}}

\begin{document}
\maketitle

\begin{abstract}
PATV-X is an explainable video-forensics repository centered on training-free temporal and frequency-domain cues for face-centric deepfake detection. The core detector operates without training on deepfake labels and is evaluated on a committed 200-video FaceForensics++ faceswap protocol. In the current artifact snapshot, the raw training-free detector reaches an AUC of 0.7325 with 0.70 balanced accuracy, while a learned gradient-boosting head trained on the extracted forensic features reaches an AUC of 0.83 with higher specificity but lower recall. The repository therefore supports a credible research story around interpretable frequency forensics and score fusion, with clear separation between the training-free core method and the optional learned extension.
\end{abstract}

\section{Introduction}
Deepfake detection systems are often accurate but opaque, or explainable but weak. PATV-X studies a middle ground: structured, interpretable forensic features combined with a transparent evaluation pipeline and optional lightweight learned heads.

\section{Repository Scope}
The repository includes:
\begin{itemize}
\item a training-free detector with multiple analysis levels,
\item feature extraction and pipeline tooling,
\item optional learned heads built on the extracted features,
\item committed evaluation artifacts under \texttt{pipeline\_results\_final/}.
\end{itemize}

\section{Method}
PATV-X combines several evidence channels, including residual motion behavior, boundary artifacts, semantic consistency checks, and multi-scale frequency forensics. The frequency layer is the most distinctive part of the repository and includes Laplacian, wavelet, SRM, and block-consistency descriptors. Final scoring is driven by an explicit weighted fusion rule rather than by an end-to-end opaque classifier.

\section{Experimental Setup}
The committed evaluation pipeline uses 200 videos from a FaceForensics++ faceswap setup (100 authentic, 100 manipulated) with pair-aware splitting to reduce identity leakage. Threshold selection is performed on validation data only, and the final report distinguishes between:
\begin{itemize}
\item the raw training-free detector,
\item an optional learned core head,
\item an open track with learned models on the extracted feature table.
\end{itemize}

\section{Current Results}
The main committed results come from \texttt{pipeline\_results\_final/pipeline\_report\_fixed.json}. The raw detector reaches AUC 0.7325, recall 0.80, specificity 0.60, and balanced accuracy 0.70 on the held-out test split. The learned core head reaches AUC 0.83, specificity 0.85, recall 0.60, and balanced accuracy 0.725.

These numbers support two restrained conclusions:
\begin{enumerate}
\item the training-free detector is credible and nontrivial,
\item the learned head improves ranking quality and specificity, but it changes the character of the method from fully training-free inference to feature-level supervised calibration.
\end{enumerate}

\section{Interpretation Boundary}
PATV-X should be presented as a serious explainable forensics repository with a validated FaceForensics++ protocol. It should not be described as already proven across all deepfake families or video domains, because the strongest committed evidence remains tied to the current dataset and split design.

\section{Conclusion}
PATV-X is technically implemented, explainable, and backed by real held-out metrics. Its strongest contribution is the repository's transparent frequency-forensics pipeline and the clear distinction it maintains between a training-free core detector and an optional learned extension.

\bibliographystyle{IEEEtran}
\bibliography{patv_x_refs}
\end{document}
